\Needspace{12\baselineskip}
A potential whose curvature is a divisorial current plus a smooth form is determined, up to ordinary domination, by its divisorial coefficients. The comparison extends to countably many components of finite total mass.

\begin{lemma}\label{lem:divisorial-potential-domination}
Let $Y$ be a connected compact complex manifold. Let $D_j$ be distinct prime divisors, $c_j>0$ real coefficients such that $S=\sum_jc_j[D_j]$ is a positive current of finite mass, and $\beta$ a smooth real $(1,1)$-form; $[D_j]$ denotes the current of integration. Let $v\in\QPSH(Y)$ satisfy
\[
 \ddc v=S-\beta.
\]
For $a\in\QPSH(Y)$, the following conditions are equivalent:
\[
 a\preceq v,\qquad a\preceq_{\mathcal I}v,\qquad
 \nu(a,D_j)\geq c_j\quad\hbox{for every }j.
\]
\end{lemma}
\begin{proof}
The coefficient-measure restriction of $S$ to $D_j$ is $c_j[D_j]$. Distinct divisors have intersections of codimension at least two, which carry no mass of either divisor current. Countable summation preserves this assertion. Smooth forms give every divisor zero coefficient mass. Local Siu decomposition after a smooth correction therefore gives $\nu(v,D_j)=c_j$.

Ordinary domination implies ideal domination by integrability, and ideal domination implies the asserted divisorial Lelong inequalities by \cref{thm:local-i-inputs}. It remains to prove that these inequalities imply ordinary domination.

Choose a Hermitian form $\Lambda$ with $a\in\PSH(Y,\Lambda)$. By \cref{lem:nonclosed-divisorial-removal}, the divisorial inequalities allow removal of each finite partial sum of $S$ from $\Lambda+\ddc a$ while retaining positivity. Passing to the increasing finite sums gives
\[
 \Lambda+\ddc a-S\geq0.
\]
Consequently $\ddc(a-v)\geq\beta-\Lambda$. The $L^1$ difference has a quasi-plurisubharmonic representative and is bounded above on compact $Y$. This yields $a\leq v+C$ almost everywhere. \Cref{lem:qpsh-almost-everywhere-order} extends this inequality everywhere. This proves ordinary domination.
\end{proof}
